% A2 v6 mathematical revision, September 9, 2026.
% Controlling report: 9975aa037d5d9eb1f339b9220f9cd21fb54c0876.
% Reviewed manuscript: 1e57d9c024f90f0304e5572c0e320707bab88074.
% This file contains new mathematical sections. It is not, by itself,
% a claim that the complete manuscript has been assembled or compiled.
% Its inherited labels refer to A2-v5-statistical-contact-rigidity.

\section{Relative locality for uniformly convex twist chains}
\label{sec:v6-locality}
The distinction between an absolute action estimate and a relative flux
estimate has a consequence which is independent of periodicity. A long
interior can change the quadratic transmission coefficient substantially
without appreciably changing the normalized nonlinear law seen at its
ends. We formulate the assertion for scalar variational chains before
specializing its interpretation to billiard length actions.

Let $L_i$, $i\in\mathbb Z$, be smooth functions on a common square,
with uniform bounds on derivatives of every fixed order. Smooth
finite-dimensional parameters are allowed, with the same uniformity.
Assume throughout the square that
\begin{equation}\label{eq:v6-chain-class}
 \begin{gathered}
 L_i(0,0)=0,\qquad DL_i(0,0)=0,\qquad
 0<b_-\le-L_{i,uv}\le b_+,\\
 L_{i,uu}\ge |L_{i,uv}|+\eta_0,\qquad
 L_{i,vv}\ge |L_{i,uv}|+\eta_0,
 \qquad \eta_0>0.
 \end{gathered}
\end{equation}
For $m<n$, put $I=(m,n)$ and $N=n-m$. Let $W_I(u,v)$ be the
stationary value of $\sum_{i=m}^{n-1}L_i(x_i,x_{i+1})$ with
$x_m=u$, $x_n=v$, and define
\[
 d_I^0=-W_{I,uv}(0,0),\qquad b_I(u,v)=-W_{I,uv}(u,v)/d_I^0.
\]
The empty interior determinant has value one. Superscripts $+$ and
$-$ denote the right half-line starting at a site and the left half-line
ending at a site, respectively. Their stationary action sums are
$S_m^+$ and $S_n^-$. Write $H_m^+(u)$ for the right interior Hessian,
$G_m^+=(H_m^+(0))^{-1}$, and $\Delta H_m^+(u)=H_m^+(u)-H_m^+(0)$.
The normalized boundary amplitude is
\begin{equation}\label{eq:v6-half-amplitude}
 B_m^+(u)=\exp\left\{
 \sum_{i=m}^\infty\log
 \frac{-L_{i,uv}(x_i^{m,+}(u),x_{i+1}^{m,+}(u))}{-L_{i,uv}(0,0)}
 -\log\det(I+G_m^+\Delta H_m^+(u))\right\}.
\end{equation}
The left amplitude is defined in the same way. The logarithmic
Fredholm determinant is defined by its trace series on the small
endpoint neighborhood constructed below.

\begin{theorem}[Relative factorization for variable chains]
\label{thm:v6-chain}
Under \eqref{eq:v6-chain-class}, the small finite stationary segments
and the small decaying half-line segments exist uniquely on one
endpoint neighborhood, uniformly in the sites and in $N$. Their
boundary actions and amplitudes satisfy
\[
 S(0)=S'(0)=0,\qquad S''\ge c_*>0,\qquad B>0,\qquad B(0)=1.
\]
There are $r_*>0$ and $0<\theta<1$ such that, for every fixed $k$,
\begin{equation}\label{eq:v6-factorization}
 \|W_I-S_m^+\oplus S_n^-\|_{C^k}
 +\|b_I-B_m^+\otimes B_n^-\|_{C^k}\le C_k\theta^N
\end{equation}
on $|u|,|v|<r_*$. The norms include the allowed parameter derivatives.
The action difference has zero value and differential at the origin,
and is bounded by $C\theta^N(|u|+|v|)^2$. The amplitude difference
vanishes at the origin. Constants at fixed orders can depend on the
corresponding higher derivative bounds. Uniform holomorphic extensions
of the edge actions give analytic constructions on a smaller common
neighborhood.
\end{theorem}
\begin{proof}
Every interior Hessian is a symmetric tridiagonal matrix with negative
off-diagonal entries and diagonal-dominance margin at least $2\eta_0$.
Write it as $D-B$, with $D$ diagonal and $B$ nonnegative. Uniform upper
bounds for $D$ imply $\|D^{-1}B\|_\infty\le q<1$. The Neumann series
therefore defines its inverse on bounded sequences. Symmetry gives the
same bound on $\ell^1$, and coercivity gives the positive inverse on
$\ell^2$. A nearest-neighbor path from $i$ to $k$ has at least
$|i-k|$ steps, whence
\begin{equation}\label{eq:v6-green}
 |G_{ik}|\le Cq^{|i-k|}.
\end{equation}
For the reference finite and right half-line inverses, paths contributing
to their difference must reach the missing right boundary and return.
Consequently
\begin{equation}\label{eq:v6-green-comparison}
 |(G_I-G_m^+)_{ik}|\le Cq^{(n-i)+(n-k)},\qquad m<i,k<n.
\end{equation}
The reflected bound holds at the other end. Differentiating the
Neumann products introduces a fixed polynomial in path length and
bounded differentiated factors; summation after a strict exponential
margin gives the corresponding fixed-order bounds. Equivalently, one
can differentiate $HG=I$ and convolve the exponentially decreasing
kernels with the bounded band matrices $\partial^\alpha H$.

Choose $q<\rho<1$. The linear Dirichlet solution is bounded by
$C(|u|\rho^{i-m}+|v|\rho^{n-i})$. The local nonlinear remainder in the
stationary gradient is quadratic. Its derivative on a ball of small
radius in the weighted maximum norm has uniformly small norm.
Convolution with \eqref{eq:v6-green} is bounded in that weighted norm.
The equation obtained by applying the reference Green operator is
therefore a contraction, on one endpoint neighborhood for every finite
interval and for both half-lines. Differentiating its fixed-point
equation leaves the same invertible operator on the highest derivative.
Locality and weighted convolution give exponential endpoint decay for
every fixed derivative. Uniqueness among all sufficiently small bounded
stationary segments follows separately by subtracting their equations
and inverting the averaged diagonally dominant Hessian. The same
argument applies to small decaying half-line segments.

The action sums and all their fixed differentiated sums converge by
endpoint decay. First variation leaves only the endpoint term. The
edge quadratic forms are bounded below by
$\eta_0(z_i^2+z_{i+1}^2)$. Applying this inequality to the derivative
of a half-line stationary segment proves a uniform positive lower bound
for $S''$. Its value and first derivative vanish at zero.

Glue the two half-line segments by addition and overwrite their
endpoints by $u,v$. The missing tails and endpoint overwrites are
exponentially small. In the interior residual the linear terms cancel;
each remaining interaction contains one left and one right endpoint
weight. Its $\ell^1$ norm, and each fixed differentiated norm, is at most
$C_k(1+N)^{a_k}\rho_1^N$, for a fixed $\rho_1<1$ slower than the
initial decay. The averaged Hessian has uniformly bounded $\ell^1$
inverse. It follows that the glued segment and the finite stationary
segment differ by this bound. Differentiated differences satisfy the
same conclusion by induction. Subtracting the action sums, estimating
the missing tails and interactions, and using this $\ell^1$ correction
proves the action assertion in \eqref{eq:v6-factorization}.

It remains to prove the relative, rather than absolute, twist estimate.
The exact corner-cofactor identity is
\[
 -W_{I,uv}(u,v)
 =\frac{\prod_{i=m}^{n-1}[-L_{i,uv}(x_i,x_{i+1})]}{\det H_I}.
\]
In particular $d_I^0>0$. With $H_I=H_I^0+\Delta H_I$, division by the
identity at zero gives
\begin{equation}\label{eq:v6-relative-log}
 \log b_I=
 \sum_{i=m}^{n-1}\log\frac{-L_{i,uv}(x_i,x_{i+1})}{-L_{i,uv}(0,0)}
 -\log\det(I+G_I\Delta H_I).
\end{equation}
Tridiagonality and endpoint localization give
\[
 \|\Delta H_I\|_1\le C(|u|+|v|),\quad
 \|\partial^\alpha\Delta H_I\|_1\le C_\alpha,\quad
 \|\partial_\xi^\alpha G_I\|_{\rm op}\le C_\alpha.
\]
The sum of absolute matrix entries bounds the trace norm, so these
bounds do not contain the number of interior sites. Only the zeroth
operator norm of $T=G_I\Delta H_I$ is required to be small, say below
$1/2$. All fixed derivatives of $T$ have bounded trace norm. Therefore
\[
 \log\det(I+T)=\sum_{s\ge1}\frac{(-1)^{s+1}}s\operatorname{tr}(T^s)
\]
and all fixed differentiated series converge. Every differentiated
product retains a trace-class factor. Derivatives of the reference
Green operator need only bounded operator norms, not trace norms.

Retain the first and last $\lfloor N/3\rfloor$ interior sites.
Removing all other entries of $\Delta H_I$, including the crossing
entries, costs exponentially small trace norm. In each retained block,
the finite stationary segment may be replaced by its own half-line
segment with the same kind of error. The identity
$\det(I+AB)=\det(I+BA)$ reduces the determinant of a block-supported
perturbation to its support. Equations \eqref{eq:v6-green} and
\eqref{eq:v6-green-comparison} compare the compressed finite inverse
with the direct sum of the two half-line compressions in operator norm,
including the off-diagonal blocks. Polynomial block sizes and fixed
parameter derivatives are absorbed by an exponential margin.
Multiplication by the uniformly trace-class perturbation preserves the
exponentially small trace-norm error.

For $\|T\|,\|T'\|\le q_1<1$, telescoping powers gives
\[
 |\operatorname{tr}(T^s-(T')^s)|
 \le s q_1^{s-1}\|T-T'\|_1.
\]
The same argument for fixed differentiated products has an additional
fixed polynomial in $s$ and a bounded number of differentiated factors.
Its series is still summable. Restoring the half-line tails costs
exponentially small trace norm. Thus the logarithmic determinant in
\eqref{eq:v6-relative-log} converges to the sum of the two half-line
logarithmic determinants. The edge sums converge by the preceding
local interaction estimates. Exponentiation proves the amplitude
assertion. A single slower $\theta<1$ absorbs every fixed polynomial;
the finitely many short intervals are covered by enlarging $C_k$.
Zero values and differentials give the stated quadratic vanishing for
the action. Uniform complex contraction and the same determinant
series prove the analytic assertion.
\end{proof}

\begin{theorem}[Relative locality under interior changes]
\label{thm:v6-relative-locality}
Let two chains satisfy \eqref{eq:v6-chain-class} with common constants.
On $[m,n]$, suppose their first $L$ and last $L$ edge actions coincide,
where $2L\le n-m$. Their intervening actions may differ arbitrarily
within this class. Then, on a common endpoint neighborhood,
\begin{equation}\label{eq:v6-locality}
 \|W_I-\widetilde W_I\|_{C^k}
 +\|b_I-\widetilde b_I\|_{C^k}\le C_k\theta^L.
\end{equation}
In the presence of parameters, agreement means agreement as parameterized
edge actions. No bound for $d_I^0/\widetilde d_I^0$ is asserted.
\end{theorem}
\begin{proof}
Apply Theorem~\ref{thm:v6-chain} to the shared interval $[m,m+L]$
with right endpoint zero. Since the right boundary action is zero at
zero and its amplitude there is one,
\[
 W_{m,m+L}(u,0)=S_m^+(u)+O_{C^k}(\theta^L),\qquad
 b_{m,m+L}(u,0)=B_m^+(u)+O_{C^k}(\theta^L).
\]
The finite functions on the left are identical for the two chains.
Their left boundary actions and amplitudes therefore differ by at most
$2C_k\theta^L$. The reflected argument applies to the last $L$ edges.
Finally apply \eqref{eq:v6-factorization} to each full interval and
subtract. Since $n-m\ge2L$ and the boundary functions have uniform
fixed-order bounds, the result is \eqref{eq:v6-locality}.
\end{proof}

Let $\mathsf H_I=D^2W_I(0,0)$. Uniform convexity and stationary
elimination give uniform positive lower and upper bounds for this
$2$ by $2$ matrix. For sufficiently small $d>0$, define
\begin{equation}\label{eq:v6-normalized-law}
 Q_I(d)=\frac{\sqrt{\det\mathsf H_I}}{\pi d^2}
       \int(d-W_I(u,v))_+b_I(u,v)\,\mathrm du\,\mathrm dv.
\end{equation}
The integration is over the common contact box; for the chosen offset
collar its sublevel is compactly contained in that box.

\begin{corollary}[Locality of normalized onset laws]
\label{cor:v6-law-locality}
The functions in \eqref{eq:v6-normalized-law} have smooth right
extensions with $Q_I(0)=1$. Under the hypotheses of
Theorem~\ref{thm:v6-relative-locality}, there is $d_*>0$ independent
of $I,L$ such that
\[
 \|Q_I-\widetilde Q_I\|_{C^k([0,d_*])}\le C_k\theta^L
\]
for every fixed $k$. For a billiard length chain with full-phase
preparation, $Q_I$ is the physical selected probability divided by
its own leading onset coefficient times $d^2$.
\end{corollary}
\begin{proof}
Write $Y=(u,v)$ and
$W_I(Y)=Y^T A_I(Y)Y/2$, where
$A_I(Y)=2\int_0^1(1-t)D^2W_I(tY)\,\mathrm dt$.
The matrices $A_I$ are uniformly positive on a smaller box. The maps
$Y\mapsto A_I(Y)^{1/2}Y$ and their inverses have uniform fixed-order
bounds there. Their differences for the two chains are bounded by
$C_k\theta^L$, since the positive matrix square root is smooth on a
compact positive matrix set and the inverse function equations can be
differentiated with uniformly bounded inverses. In these coordinates
the sublevel is the common disk $|w|^2/2<d$. The transformed amplitudes
and their fixed derivatives differ by $C_k\theta^L$.

Radial Taylor expansion on this disk cancels odd homogeneous terms.
Integration against $(d-|w|^2/2)_+$, followed by division by $d^2$,
therefore gives a smooth right function. Taylor remainders to sufficiently
high order give the claimed estimate through every fixed $d$ derivative,
including at zero. The constant term is one because
$\int(d-Y^T\mathsf H_IY/2)_+\,\mathrm dY
=\pi d^2/\sqrt{\det\mathsf H_I}$ and $b_I(0)=1$.
For billiard length actions, the physical density is
$d_I^0b_I(u,v)\,\mathrm du\,\mathrm dv\,\mathrm dr/(2\pi A)$.
Its leading probability is
$d_I^0d^2/(2A\sqrt{\det\mathsf H_I})$.
This proves the asserted normalization.
\end{proof}

The conclusion concerns a normalized nonlinear law. It does not compare
raw probabilities after an unknown change of their exponentially small
leading factors. Nor does it assert that every abstract chain is realized
by a billiard configuration. For alternating length actions the local
hypotheses follow from positive facing curvatures, and geometric
localization and the full-phase measure provide the additional physical
identification. This separates the transferable relative determinant
mechanism from its billiard realization.

\section{Coefficient sensitivity and the contact inverse}
\label{sec:v6-sensitivity}
The equality of the information determined by two exact analytic germs
need not imply comparable recovery from noisy observations. The
following comparison places the one-flight and limiting contact-jet
formulas in the same coefficient coordinates. It records the refinement
identified in the independent referee report.

\begin{proposition}[Comparison of the top-jet diagonals]
\label{prop:v6-sensitivity}
In the identical-even contact class, put
$D_m^\infty=|\partial_{q_{2m}}f_{m-1}|$ and
$D_m^{(1)}=|\partial_{q_{2m}}f_{1,m-1}|$. Then
\begin{equation}\label{eq:v6-sensitivity-ratio}
 \frac{D_m^\infty}{D_m^{(1)}}
 =(\tanh\gamma)^m\left[\coth(m\gamma)+2mD_m\right],\qquad
 D_m=\frac1{e^{2(m-1)\gamma}-1}-\frac1{e^{2m\gamma}-1}.
\end{equation}
For fixed $\gamma>0$ the ratio is
$(\tanh\gamma)^m[1+O_\gamma(m e^{-2(m-1)\gamma})]$ as $m\to\infty$.
\end{proposition}
\begin{proof}
The formulas \eqref{eq:v5-diagonal} and \eqref{eq:v5-one-flight}
have the same factor $4/[2^m(m+1)(m!)^2]$. The remaining factor in
the quotient is $(a\nu_1)^{-m}$, where
$a=\sinh\gamma/g$ and $\nu_1=g\cosh\gamma/\sinh^2\gamma$.
Thus $a\nu_1=\coth\gamma$, proving
\eqref{eq:v6-sensitivity-ratio}. The asserted remainder follows by
expanding $\coth(m\gamma)-1$ and bounding the two positive exponential
denominators defining $D_m$.
\end{proof}

This proposition compares the differential with respect to the last jet
while lower jets are fixed, at equal absolute error in the corresponding
coefficient. A common rescaling of that geometric jet cancels from the
ratio. The proposition is not a bound for the condition number of the
whole triangular inverse, for extracting coefficients from noisy
probabilities, or for analytic continuation. It does not compare minimax
risks of experiments with different sampling costs. The exact germ
uniqueness and the fixed finite-jet inverse remain as stated in
Theorem~\ref{thm:v5-jets}; the relative long-chain theorem supplies the
uniform limiting physical law, not geometric parameters absent from the
one-flight nonlinear germ.

\section{Three observations in the physical family}
\label{sec:v6-three-observations}
The scalar model \eqref{eq:v5-amplitude-model} concerns actual contact
curvatures when the two facing curvatures in each channel are equal.
Without that assumption its nodes are effective channel parameters:
$c_i=\sqrt{(1+g\kappa_{i,0})(1+g\kappa_{i,1})}$, not separate
endpoint curvatures. The general forward theorem is unchanged.

The four-amplitude inverse allows area to vary independently. In the
physical support family \eqref{eq:g-isogap-family}, area is instead a
specified symmetric function of the three radius deviations. Let
$x_i=r_i-R$ and let $e_1,e_2,e_3$ be their elementary symmetric
polynomials. Direct substitution in \eqref{eq:v3-area} gives
\begin{equation}\label{eq:v6-physical-area}
 A_R(e)=A_0-\frac{\pi R}{54}e_1
             +\frac{41\pi}{7776}e_1^2-\frac{5\pi}{432}e_2,
 \qquad A_0=\frac{\sqrt3}{2}-\pi R^2,
 \qquad g=1-2R.
\end{equation}
Indeed, $\alpha=e_1/108$ and
$\beta^2+\zeta^2=e_1^2/2916-e_2/972$. These identities also show
that the constrained area introduces no additional independent unknown.

\begin{theorem}[Three-amplitude local recovery]
\label{thm:v6-three}
For every fixed $R\in(0,1/2)$, there is a set
$J_R\subset\{1,2,3,4\}$ of cardinality three such that the amplitudes
$(C_j)_{j\in J_R}$ determine the unordered curvature triple and the
area locally near the physical circle with radius $R$. For two members
in a sufficiently small neighborhood at the same $R$,
\[
 |A-\widetilde A|\le C\delta_{J_R},\qquad
 \operatorname{dist}_{\rm match}(\kappa,\widetilde\kappa)
       \le C\delta_{J_R}^{1/3},\qquad
 \delta_{J_R}=\max_{j\in J_R}|C_j-\widetilde C_j|.
\]
The estimates preserve multiplicities. At $R=1/4$, one may take
$J_R=\{1,2,3\}$. The exponent $1/3$ is optimal. A compact interval
of reference radii is covered by finitely many such observation charts.
If the gap is also measured, the estimates on one chart hold with
$\delta_{J_R}+|g-\widetilde g|$ on the right.
\end{theorem}
\begin{proof}
Use the analytic coefficient extension $M_j(e)$ from
\eqref{eq:v5-contour}. The ambient four-amplitude map
$(A,e)\mapsto(A^{-1}M_j(e))_{j=1}^4$ has invertible differential at
$(A_0,0)$, by \eqref{eq:v5-jacobian} and
\eqref{eq:v5-wronskian}. The graph embedding
$e\mapsto(A_R(e),e)$ has rank three. Its composition with that
ambient map therefore has rank three. Some three rows have nonzero
determinant. The analytic inverse function theorem for those rows,
applied in the full coefficient neighborhood before restricting to
real-rooted cubics, proves locally Lipschitz recovery of $e$.
Equation~\eqref{eq:v6-physical-area} gives the area estimate. The
multiplicity-preserving cubic root estimate in the proof of
Theorem~\ref{thm:v5-pairwise}, followed by the Lipschitz map $r\mapsto1/r$
on a positive compact interval, gives the curvature estimate.

We record explicitly the first-three calculation at $R=1/4$, as
identified in the independent report. At $e=0$ the rows of the
three-dimensional differential are
\[
 \frac1{A_0}\left(
 \Phi_j'(R)+\frac{\pi R}{18A_0}\Phi_j(R),\quad
 -\Phi_j''(R)+\frac{5\pi}{144A_0}\Phi_j(R),\quad
 \frac12\Phi_j'''(R)\right).
\]
At $g=1/2$, direct differentiation gives the following values divided
by $\sqrt2$:
\[
\begin{array}{c|rrrr}
 j&\Phi_j&\Phi_j'&\Phi_j''&\Phi_j'''\\\hline
 1&1/4&3/4&-5/4&21/4\\
 2&1/24&17/72&35/216&-491/216\\
 3&1/140&297/4900&36243/171500&-4458537/6002500
\end{array}
\]
Substitution yields
\begin{equation}\label{eq:v6-three-jacobian}
 \det D\Psi_3(0)
 =-\frac{2\sqrt2(15804720A_0+64253\pi)}{72930375A_0^4}\ne0,
 \qquad A_0=\frac{\sqrt3}{2}-\frac\pi{16}>0.
\end{equation}
Its sign is exact, not inferred from a floating-point determinant.
The same minor stays nonzero on an interval about this reference.

For general $R$, continuity of a chosen nonzero minor supplies a local
chart. Compactness gives a finite cover of a compact reference interval,
with a possibly different choice of three amplitudes on each chart.
This does not assert that the first three work at every radius. If $R$
is variable, use the data $(g,C_j:j\in J_R)$ and coordinates $(R,e)$.
Since $g=1-2R$, their differential is block triangular with nonzero
diagonal blocks. The inverse function theorem gives the final assertion.

The physical two-sided path of Theorem~\ref{thm:v5-sharpness} has
matching curvature distance comparable to $|s|$ while every fixed
selected amplitude difference is $O(|s|^3)$. An exponent larger than
$1/3$ would contradict these bounds. This also proves sharpness for
any of the selected observation charts without requiring a separately
nonzero cubic coefficient in each selected amplitude.
\end{proof}

The rank argument uses the four-amplitude theorem to select a sufficient
three-dimensional physical observation. It does not remove the fourth
observation from the model with independently unknown area. Both
statements have different parameter spaces, and both are retained.

\section{Count-only acquisition with charged onset localization}
\label{sec:v6-count-acquisition}
A leading-amplitude inverse and a conditional-position inverse are not
the same experiment. We now connect the former to raw count records.
All preparations in this section are independent full-phase preparations;
no independence between successive impacts is used. Observation times
and collision counts are recorded without error. The constants depend
on a fixed extrapolation order and on the specified compact family.

Fix a sufficiently small compact physical family in an observation chart
of Theorem~\ref{thm:v6-three}, and let $J$ be its selected three flight
numbers. The same construction works with all four amplitudes.
For $j\in\{1,2,3,4\}$, write
\begin{equation}\label{eq:v6-count-expansion}
 p_j(d)=\Pr(N_{jg+d}\ge j+1)=d^2F_j^{\rm amp}(d),\qquad
 F_j^{\rm amp}(0)=C_j.
\end{equation}
On a common collar these functions have uniformly bounded fixed-order
right derivatives, and are bounded above and below by positive constants.
These are the unlabelled physical probabilities, not selected endpoint
observations. The threshold theorem supplies
\eqref{eq:v6-count-expansion} for this equal-gap family.

\begin{lemma}[Charged coarse localization]
\label{lem:v6-coarse-search}
Suppose a fixed interval $[a_0,b_0]$ containing the gap is supplied,
with positive endpoints and width at most a common one-flight onset
collar. Its width does not depend on the requested accuracy. For every
sufficiently small $h>0$ and $0<\eta<1/2$, a count-only experiment
using at most $C h^{-2}\log(C/\eta)$ preparations returns an interval
of width at most $h/16$ containing the gap with probability at least
$1-\eta$. The cost bound is deterministic, including on unsuccessful
outcomes.
\end{lemma}
\begin{proof}
For a current interval $[a,b]$ of width $w>h/16$, set $t=(a+b)/2$.
At time $t$, perform
\[
 n(w)=\left\lceil Cw^{-2}
       \left[\log(C/\eta)+\log(32w/h)\right]\right\rceil
\]
fresh first-flight trials and record whether $N_t\ge2$. If at least
one succeeds, replace the interval by $[a,t]$. If none succeeds,
replace it by $[t-w/4,b]$. A success implies $g<t$ exactly, since the
count event is impossible at and below the onset. Conditional on the
current interval containing $g$, the no-success update can exclude $g$
only when $t-g>w/4$. In that case
$p_1(t-g)\ge c w^2$, by \eqref{eq:v6-count-expansion}; hence the
conditional failure probability is at most $c'\eta h/w$, with $c'$
as small as needed by the choice of $C$.

Each update reduces the width by either $1/2$ or $3/4$, and the
algorithm stops when it is at most $h/16$. Along every path, the widths
before stopping satisfy $\sum h/w\le C_0$. This follows by reading
the sequence backwards from its last width and using a geometric ratio
at least $4/3$. The sum of conditional failure probabilities before the
first invalid bracket is therefore at most $\eta$. The same backwards
geometric estimate, with its extra logarithmic factor, gives
\[
 \sum n(w)\le C h^{-2}\log(C/\eta).
\]
The ceiling contributions are harmless because the number of steps is
logarithmic in $1/h$. These estimates hold for every sequence of
outcomes, not just the successful ones.
\end{proof}

\begin{theorem}[Charged count-only recovery near coalescence]
\label{thm:v6-count-acquisition}
Fix an integer $m\ge1$. Under the preceding physical and observation
assumptions, including the fixed initial gap interval of
Lemma~\ref{lem:v6-coarse-search}, there is an experiment based only
on counts which produces $\widehat g$, $\widehat A$, and an unordered
curvature triple $\widehat\kappa$. With probability at least $1-\eta$,
\begin{equation}\label{eq:v6-count-errors}
 |\widehat g-g|\le C_m h^{m+1},\qquad
 |\widehat A-A|\le C_m h^m,\qquad
 \operatorname{dist}_{\rm match}(\widehat\kappa,\kappa)
                 \le C_m h^{m/3}.
\end{equation}
The total number of preparations, including coarse localization, fine
calibration, and amplitude estimation, is at most
\begin{equation}\label{eq:v6-count-cost}
 T\le C_m h^{-(2m+2)}\log(C_m/\eta)
\end{equation}
on every outcome. In particular a sufficient count-only order for
curvature accuracy $\varepsilon$ is
\begin{equation}\label{eq:v6-count-epsilon}
 C_m\varepsilon^{-(6+6/m)}\log(C_m/\eta).
\end{equation}
The assertion is a sufficient upper bound, not a minimax statement.
\end{theorem}
\begin{proof}
Allocate the failure probability among the three stages. Run
Lemma~\ref{lem:v6-coarse-search} with a suitable fraction of $\eta$
and let $g_0$ be the midpoint of its output. On its success event,
$|g-g_0|\le h/32$. In particular every one-flight programmed offset
$g_0+lh-g$, $1\le l\le m+1$, is positive and comparable to $h$.

At each of those $m+1$ times take a deterministic number
\[
 n=\left\lceil C_mh^{-(2m+2)}\log(C_m/\eta)\right\rceil
\]
of fresh preparations. Let $\widehat p_l$ be the empirical probability.
Conditional on a successful coarse stage, the true probabilities are
bounded above and below by fixed multiples of $h^2$. The binomial
exponential bound therefore gives relative errors at most $c_mh^m$
simultaneously, outside a further event of probability at most a fixed
fraction of $\eta$. Apply the square-root interpolation procedure of
Lemma~\ref{lem:v5-calibration} to these unlabelled one-flight
probabilities. Its proof uses only a factorization into an unknown
positive coefficient times $d^2$ times a smooth function equal to one
at zero; it therefore applies to \eqref{eq:v6-count-expansion} with
coefficient $C_1$ as well. Let $\widehat\tau$ be the unique root in
$[-h/2,h/2]$ when it exists, and zero otherwise. Then
$\widehat g=g_0+\widehat\tau$ obeys
$|\widehat g-g|\le C_mh^{m+1}$ on the simultaneous successful event.
No negative-offset physical probability is used in this root argument.

For amplitude estimation use fresh preparations at times
$j\widehat g+lh$, $j\in J$, $1\le l\le m$, with $n$ trials per time.
Put $\omega_l=(-1)^{l-1}\binom ml$ and define
\[
 \widehat C_j=\sum_{l=1}^m\omega_l
 \frac{\widehat p_{j,l}}{(lh)^2}.
\]
On the fine-calibration success event, the actual offset is
$d_{j,l}=lh+j(\widehat g-g)$, which is between fixed positive
multiples of $h$ for sufficiently small $h$. Since $j\le4$, smoothness
in \eqref{eq:v6-count-expansion} gives
\[
 \left|\frac{p_j(d_{j,l})}{(lh)^2}-F_j^{\rm amp}(lh)\right|
       \le C_m\frac{|\widehat g-g|}{h}\le C_mh^m.
\]
The identities $\sum_l\omega_l=1$ and
$\sum_l\omega_l l^a=0$ for $1\le a<m$ cancel the lower Taylor terms
of $F_j^{\rm amp}$. Thus the population extrapolation differs from
$C_j$ by at most $C_mh^m$.

For completeness, at a positive offset comparable to $h$, Bernstein's
inequality for the empirical Bernoulli mean gives
\[
 \left|\frac{\widehat p-p}{(lh)^2}\right|
 \le C_m\left[
 \sqrt{\frac{\log(C_m/\eta)}{nh^2}}
 +\frac{\log(C_m/\eta)}{nh^2}\right]
\]
with the requisite confidence. This follows from $p\le C_mh^2$,
not from treating the normalized Bernoulli variable as having bounded
variance independent of $h$. Conditional on the complete pilot, the
new trials remain independent. The selected value of $n$ and a union
bound over the fixed number of times therefore give
$\max_{j\in J}|\widehat C_j-C_j|\le C_mh^m$.

Use the local inverse for the data $(g,C_j:j\in J)$ in
Theorem~\ref{thm:v6-three}. With noisy data, minimize their maximum
discrepancy over a small compact physical coefficient set containing
the true member in its interior relative to the physical set. A
minimizer exists. Since the true member has discrepancy at most
$C_mh^m$, the fitted member's exact data differ from the true data by
at most $2C_mh^m$. The locally Lipschitz coefficient inverse and the
multiplicity-preserving root estimate give the last two bounds in
\eqref{eq:v6-count-errors}. The gap estimator in that display is the
fine pilot estimator, not a less accurate gap coordinate of a joint
least-discrepancy fit.

Every trial count after coarse localization is deterministic. On a
failed pilot, a programmed time may have a nonpositive true excess;
the experiment still performs only the specified finite number of
trials. Thus no uncontrolled waiting time is introduced on failure
outcomes. Lemma~\ref{lem:v6-coarse-search} also has a deterministic
cost bound. Summing its cost and the fixed number of batches proves
\eqref{eq:v6-count-cost}. The confidence follows from the successive
conditional bounds and a union bound. Finally choose
$h=c_m\varepsilon^{3/m}$, with $c_m$ small enough for the curvature
constant in \eqref{eq:v6-count-errors}. This proves
\eqref{eq:v6-count-epsilon}.
\end{proof}

If the gap is known, the two localization stages can be omitted, with
the same sufficient exponent. The charged theorem still assumes a fixed
prior compact family and a fixed initial interval lying inside a common
onset collar; it does not conceal an accuracy-dependent supplied bracket.
It is not a global search over arbitrary unknown billiard tables.

The observations and conclusions can now be kept distinct. Exact
normalized probability germs determine the analytic even contact in the
specified class. Finite coefficient vectors give the fixed finite-jet
inverse, but their extraction from noisy probabilities is a further
problem. Unlabelled leading amplitudes determine a curvature multiset
with a cubic-root modulus near coalescence; the count-only experiment
above supplies a corresponding preparation upper bound. Selected
conditional positions retain a different observable and satisfy the
separate acquisition theorem with order
$\varepsilon^{-(2+2/m)}$ at fixed flight order, including its stated
fine-calibration stage. No preparation rate is transferred from one of
these observations to another merely because both involve collisions.
